\subsection{平展代数的子代数}

命题3. --- 设A是K上的平展代数。A的子代数和理想只有有限多个。此外，K的每一个使A对角化的扩张也使A的每个子代数和每个商代数对角化；特别地，这些代数都是平展的。

只需证明代数 \(K^{n}\) 只有有限多个子代数和理想，并且 \(K^{n}\) 的子代数和商代数都是可对角化的。我们用 \((\varepsilon_{1},\ldots,\varepsilon_{n})\) 表示 \(K^{n}\) 的典范基。

设 A 是 \(K^{n}\) 的一个子代数，并令 \(v_{1}, \ldots, v_{n}\) 为 n 个投影 \(K^{n} \rightarrow K\) 在 A 上的限制。由于这些 \(v_{i}\) 的核之交显然为 0，这些 \(v_{i}\) 生成与 A 对偶的向量 K-空间（II，第302页，推论1）；因此 K-代数 A 是可对角化的（V，第29页，命题1）。

对 \(\{1, 2, \ldots, n\}\) 的每个子集 I，令 \(\varepsilon_{I} = \sum_{i \in I} \varepsilon_{i}\) 。显然，元素

\(\varepsilon_{I}\) 是 \(K^{n}\) 的幂等元；我们有 \(\varepsilon_{I}=0\) 当且仅当 I 为空，且 \(\varepsilon_{I}\varepsilon_{J}=\varepsilon_{I\cap J}\) 。根据前述讨论，\(K^{n}\) 的每个子代数 A 都是可对角化的；因此根据命题1的条件 b)，\(K^{n}\) 的每个子代数 A 都具有一个基 \((\varepsilon_{I_{1}},\ldots,\varepsilon_{I_{p}})\) ，其中 \((I_{1},\ldots,I_{p})\) 是 \(\{1,2,\ldots,n\}\) 的一个划分，而这样的子代数只有有限多个。

对 \(\{1,2,\ldots,n\}\) 的每个子集 I，令 \(a_{I}\) 为 \(K^{n}\) 的向量子空间，它以幂等元 \(\varepsilon_{i}\) （对 \(i\in I\)）为基；那么显然 \(a_{I}\) 是 \(K^{n}\) 的一个理想；此外，若 \(J=\{1,2,\ldots,n\}-I\) ，则剩余类 \(\overline{\varepsilon}_{j}\) （即 \(\varepsilon_{j} \mod a_{I}\) ，对 \(j\in J\)）构成 \(K^{n}/a_{I}\) 的一个基。我们有 \(\overline{\varepsilon}_{j}^{2}=\overline{\varepsilon}_{j}\) 且 \(\overline{\varepsilon}_{j}\overline{\varepsilon}_{k}=0\) （若 \(j\neq k\)），因此代数 \(K^{n}/a_{I}\) 是可对角化的，由 V 第29页命题1。

仍需证明 \(K^{n}\) 的每个理想都具有形式 \(a_{I}\) 。设 I 为满足 \(1 \leqslant i \leqslant n\) 且 \(\varepsilon_{i} \in a\) 的整数 i 的集合，则 \(a_{I} \subset a\) 。设 \(x = x_{1}\varepsilon_{1} + \cdots + x_{n}\varepsilon_{n}\) 是 a 的一个元素（其中 \(x_{1}, \ldots, x_{n}\) 在 K 中），并设 i 属于 \(\{1, 2, \ldots, n\} - I\) 。我们有 \(x_{i}\varepsilon_{i} = x\varepsilon_{i} \in a\) ，以及 \(\varepsilon_{i} \notin a\) ，从而 \(x_{i} = 0\) 。因此 \(x = \sum_{i \in I} x_{i}\varepsilon_{i}\) ，而这表明

\(x \in a_{1}\) 。我们现在已经证明了 \(a \subset a_{1}\) ，由此最终 \(a = a_{1}\) 。

推论。——设 \(A_{1}, \ldots, A_{m}\) 是K上的代数，且 \(A = A_{1} \times \cdots \times A_{m}\) 。要使 A 为平展的，必要且充分条件是 \(A_{1}, \ldots, A_{m}\) 是平展的。

假设A是平展的；每个代数 \(A_{1}, \ldots, A_{m}\) 都同构于A的一个商，因此由命题3可知它是平展的。反之，K 的每一个使 \(A_{1}, \ldots, A_{m}\) 对角化的扩张显然也使 A 对角化。

